\chapter{Palnteamiento del problema}
La creciente generación de residuos orgánicos provenientes de actividades agroindustriales y urbanas representa un reto ambiental urgente a nivel global. Su acumulación descontrolada no solo degrada el suelo, el agua y el aire, sino que también constituye una fuente significativa de gases de efecto invernadero (GEI), particularmente dióxido de carbono (CO$_2$) y metano (CH$_4$), responsables del agravamiento del cambio climático \parencite{FAO2021,IPCC2014}. Si bien existen tecnologías tradicionales como el compostaje o la digestión anaeróbica, estas presentan limitaciones económicas, operativas y ambientales, sobre todo en contextos rurales o de escasos recursos \parencite{Smith2014}. Este panorama exige el desarrollo de soluciones más eficaces, sostenibles y adaptables que permitan valorizar los residuos y mitigar sus impactos.

Una de las alternativas más prometedoras en este contexto es la bioconversión de residuos orgánicos mediante larvas de \textit{Hermetia illucens}. Esta especie ha demostrado ser eficaz en la transformación de materia orgánica en productos útiles como proteína y larvicompost, contribuyendo a reducir la presión sobre vertederos y a cerrar ciclos de nutrientes. Sin embargo, el conocimiento sobre las emisiones de GEI generadas durante este proceso aún es fragmentario \parencite{bermudezComprehensiveUtilizationBlack2023,Surendra2016}. Muchos estudios ofrecen solo estimaciones estáticas, sin considerar variables ambientales clave como la temperatura, la humedad o el tipo de sustrato, lo cual limita el desarrollo de estrategias de mitigación basadas en evidencia \parencite{fuhrmannComprehensiveIndustryrelevantBlack2025,rossiEstimatingDynamicsGreenhouse2024}.

A esto se suma una débil integración de tecnologías emergentes, como sensores del Internet de las Cosas (IoT) y algoritmos de aprendizaje automático (AAA), en los sistemas de bioconversión. La falta de soluciones que combinen monitoreo en tiempo real, modelado matemático y predicción adaptativa de emisiones, obstaculiza la transición hacia procesos verdaderamente inteligentes, sostenibles y trazables. Sin sistemas robustos de medición, reporte y verificación (MRV), resulta difícil evaluar, replicar y optimizar estos procesos a distintas escalas y condiciones operativas \parencite{kamilaris2017review}.

En respuesta a estas brechas, esta investigación propone el desarrollo de un sistema predictivo híbrido que integre sensores IoT de bajo costo, algoritmos de AAA y modelos matemáticos basados en ecuaciones diferenciales. Este sistema permitirá cuantificar y predecir en tiempo real las emisiones de CO$_2$ y CH$_4$ durante la bioconversión con \textit{Hermetia illucens}, considerando variables como el tipo de residuo, la humedad del sustrato, la temperatura y la composición larval. Con ello, se busca avanzar hacia una gestión de residuos más eficiente, ambientalmente responsable y alineada con los principios de economía circular y mitigación climática \parencite{barragan2017nutritional,Diener2011,salamEffectDifferentEnvironmental2022}.

